\question[10]
Explain the concept and goals of maximum sustainable yield (MSY) as
it relates to commercial fisheries. What population / life history
parameters are required to estimate MSY? Why is MSY hard to implement in
practice?

\begin{solution}

The goal of MSY is to sustain a viable population that maximizes the harvestable yield. (2 pts)

Must consider how often the organisms reproduce, their reproductive output, maximum life span. Be rather generous here, as long as it is a reasonable population / life history parameter. (3 pts)

It's hard to implement because can't usually know or accurately estimate those parameters. (3 pts)

\end{solution}

\question[5]
Explain one reason a conservation manager must understand how genetic variation is
distributed across the entire range of a species that is being
considered for conservation management or reintroduction into a
particular area.

\begin{AnswerPage}{0.2\textheight}
	Many studies have shown that many organisms are not genetically uniform across its entire range.  
	
	\begin{itemize}
		
		\item A population adapted for one area my not be well adapted for another area. Thus, individuals are not amenable to transplant.
		
		\item Cannot let individuals in one area go extinct because they are genetically different. 
		
		\item Genetic variability across the population increases adaptive ability in a changing environment.
	\end{itemize}
	
	Grade as 4 points per explained reason.
	
\end{AnswerPage}


\question[10]
Briefly explain each of the three factors that must be considered to identify
a potential biodiversity hotspot for conservation efforts.

%\textbf{Can turn into individual questions. E.g., Explain why FACTOR X is important to
%	identify a potential biodiversity hotspot.}

\begin{AnswerPage}{4\basespace}%
	
	Grade as 1 point per item, plus 2 points for the explanation.\vspace{\baselineskip}
	
	1. Identify areas that have the highest organismal diversity.
	
	2. Identify the areas that have the highest amount of endemism.
	
	3. Identify the areas that have the greatest threats from human impacts.
	
	The areas where these three items intersect are the best candidates for biodiversity hotspots.
\end{AnswerPage}

%\vfill
%
%\bonusquestion
Select whether you want \textbf{2 points} or \textbf{5 points} added to your score for this exam. Before you decide, there is a small catch: if more than 10\% of the class selects 5 points, then no one gets the extra points. Only I will see your choice.

\begin{oneparchoices}
	\choice 2 points.
	\choice 5 points.
\end{oneparchoices}
	
%
